\chapter{Εγκατάσταση και αναπαραγωγή}
\label{app:build}

\section{Εφαρμογή συσκευής}

\begin{enumerate}
  \item Unity 6 (6000.5) με υποστήριξη Android. Άνοιγμα του φακέλου \code{QuestMarkerApp/} από το Unity Hub.
  \item Το Meta XR Core SDK και τα υπόλοιπα πακέτα κατεβαίνουν αυτόματα από το \code{Packages/manifest.json}.
  \item Από το μενού \code{Quest Marker App} του Editor εφαρμόζονται οι ρυθμίσεις build για Quest (ελάχιστο API 32, στόχος API 34, IL2CPP, ARM64).
  \item Σύνδεση του Quest~3 με USB σε λειτουργία προγραμματιστή και \code{Build and Run}.
  \item Προαιρετικά, για φωνητικές εντολές: πακέτο \code{com.meta.xr.sdk.voice}, ρύθμιση Wit.ai και σύμβολο \code{META_VOICE_SDK}.
\end{enumerate}

\section{Εργαλεία υπολογιστή}

\begin{lstlisting}[language=bash,caption={Εγκατάσταση και έλεγχος των εργαλείων Python.}]
cd QuestMarkerApp/Server
python -m venv .venv
.venv\Scripts\activate          # Windows
pip install -r requirements.txt
pip install pymeshlab pillow pytest   # προετοιμασία μοντέλων, tests
python -m pytest -q                   # 127 tests
\end{lstlisting}

\section{Αναπαραγωγή των πειραμάτων}

\begin{enumerate}
  \item \code{git checkout thesis-v2} για την έκδοση της εφαρμογής που χρησιμοποιήθηκε στα πειράματα.
  \item Τα ακατέργαστα δεδομένα βρίσκονται στο \code{thesis/data/}.
  \item Τα σενάρια ανάλυσης στο \code{thesis/analysis/} (\code{pip install -r requirements.txt}) γράφουν πίνακες και γραφήματα στο \code{thesis/figures/generated/}.
  \item Η αναφορά μεταγλωττίζεται με \code{latexmk} (XeLaTeX + Biber) στον φάκελο \code{thesis/report/}.
\end{enumerate}
